\chapter{Raccordement matériel entre la structure unique et ses cinq résolutions}
\label{chap:pdf2-generacion-simultanea-hmt}

\begingroup
\footnotesize
% CORTE_2_NO_DISGREGACION — bloque de empalme PDF1 -> PDF2
% Congelado después del PDF1 compilado el 2026-08-19.
% Owners:
% 04b  8ba02f246c18a3bd2ff83ffdb08874c816a7f59c688b6fab6af6eb2f9647071f
% 04b1 25d6dee3875cbdd3252f793212448c17c4b5b2d3c5d787e995598d18fa76d5cf
% 04b2 d5de9b5379d4f95ea3d35e63c9ccc2ab741b5c750c01409af6f13526073b874f
% 04b3 ac9ca321b610e2c0560a32ec6ae1555814c67c587e152096d0c48d6abaa0bd4b

\section*{Raccordement mathématique concret PDF1--PDF2}

L'entrée unique est l'état APP--TRIT--TPK complet. Avant toute évaluation archimédienne sont parcourus
\[
 0\longrightarrow10,\qquad 0\longrightarrow1000,
 \qquad 0\longrightarrow\infty,
\]
avec les neuf espaces internes \(1,\ldots,9\) ; deux positions produisent APP, TRIT oriente ses feuillets et TPK prolonge extension, phase, retenue, mémoire, frontière, parcours et survie. Ni le continu ni \(\mathbb R\) ne sont des entrées.

Pour une émission survivante
\[
 \mathfrak e_\alpha=
 (\alpha,\operatorname{Ext}_{\gamma_\alpha},\operatorname{or}(\alpha),
 \kappa(\alpha),\partial\alpha,\gamma_\alpha,
 m(\gamma_\alpha;m_0),\operatorname{Led}(\gamma_\alpha))
\]
les cinq caractéristiques sont des actions simultanées sur ce même générateur :
\begin{align*}
 \mathcal O_{\rm sp}(\mathfrak e_\alpha)
  &=(U_\alpha,\sigma_x,\sigma_y,A_\alpha,\eta_\alpha),\\
 A_\alpha&=P_yU_\alpha P_x+Q_yU_\alpha Q_x,
 &\eta_\alpha&=Q_yU_\alpha P_x+P_yU_\alpha Q_x,\\
 \mathcal O_{\rm sol}(\mathfrak e_\alpha)
  &=(b_\alpha,b_\alpha^+,b_\alpha^-;
  M^\pm\mapsto M^\pm+b_\alpha^\pm,
  D\mapsto D+b_\alpha,H\mapsto H+|b_\alpha|),\\
 \mathcal O_{\rm gau}(\mathfrak e_\alpha)
  &=\{(\operatorname{Cell}_{\rm TPK}(\alpha),q,B,W_c(q)):
      q\in Q_{\rm inc}(t\alpha)\},\\
 \mathcal O_{\rm coh}(\mathfrak e_\alpha)
  &=(I_y,J_y,\kappa_{y,N},\delta_y,C_{y,N},
      F_\alpha:C_{y,N}\to C_{x,N}),\\
 \mathcal O_{\rm det}(\mathfrak e_\alpha)
  &=(\Psi_\alpha,K_\alpha,\mathbf1_\alpha,z_\pm(\alpha),h_*(\alpha),
  \operatorname{Det}C_{y,N},H_*C_{y,N},\operatorname{Smith}C_{y,N},
  \operatorname{Boc}C_{y,N}).
\end{align*}
Ces expressions sont des vues dépendantes d'une seule action \(\mathcal O(\mathfrak e_\alpha)\) ; elles ne définissent pas cinq domaines, facteurs ou branches. Elles partagent littéralement source, cible, parcours, registre, orientation, retenue, mémoire et frontière. Leurs identités de couplage sont
\[
 A_\alpha^*A_\alpha+\eta_\alpha^*\eta_\alpha=U_\alpha^*U_\alpha,
 \qquad E b_f^\pm=b_c^\pm+\kappa^{\rm can},
\]
\[
 Q_{\rm inc}=C^2(\operatorname{Cell}_{\rm TPK};\mathbf F_3),
 \qquad sB=0,\qquad W_c(q+Ba)=W_c(q),
\]
\[
 \partial h_*=z_--z_+,qquad
 K_{\beta\alpha}=K_\beta+\Psi_\beta K_\alpha.
\]

La monodromie est la connexion nonadique réelle :
\[
 \Gamma_{9,k}=\operatorname{Hol}_{C_9}(\nabla_{\rm TPK}),
 \qquad g(k+9)=g(k),
 \qquad q_{\rm mem}(k+9)=q_{\rm mem}(k)+1,
\]
sans retour d'état. Pour une histoire \(h=(\varepsilon_{k+1},\ldots,\varepsilon_N)\), les treize coordonnées sont engendrées par le pliage
\[
 r_{k,N}(h)=
 \operatorname{Upd}_{\varepsilon_N}\circ\cdots\circ
 \operatorname{Upd}_{\varepsilon_{k+1}}(r^0_{k}),
\]
et le terminal renvoie la multisection entière, sans sélecteur :
\[
 \mathfrak T_{k,N}(\widehat x_k)
 =\{r_{k,N}(h):h\in\mathscr H_{k,N}(\widehat x_k)\}.
\]
La naturalité est une égalité unique sur les préfixes engendrés,
\[
 u_{\ell,k}^{\mathcal O}\mathcal O_{\le\ell}(\gamma)
 =\mathcal O_{\le k}(\gamma)
 =\operatorname{map}(\mathcal O)
   \bigl(u_{\ell,k}^{\mathfrak e}E_{\le\ell}(\gamma)\bigr),
\]
qui passe au système inverse des terminaux. Ce n'est qu'alors que l'on obtient la structure discrète complète du continu ; l'évaluation dans \(\mathbb R\) est une expression ultérieure.

\subsection*{Carte d'utilisation dans PDF2}

\begin{center}
\small
\begin{tabularx}{\textwidth}{@{}p{.18\textwidth}X p{.30\textwidth}@{}}
\toprule
Symbole utilisé & Contenu structurel intégral de PDF1 &
Entrelaceur construit dans PDF2 \\
\midrule
\(\mathfrak G_{\rm sp}\) &
\(U,\sigma,P,Q,A,\eta\), feuillets, retenue, terminal, naturalité et lecteur signé interne. &
\(\Sigma_R^{\rm sp},J_\pm,K_R^{\rm sp},\Xi,P,B_W\) et l'identification avec la forme de Weil. \\
\(\mathfrak G_{\rm sol}\) &
\(b,b^\pm,\kappa^{\rm can}\), mémoire, composantes micro/coquille orthogonales, colonne terminale, Gram, Stein et télescopie sans double comptage. &
Application construite depuis la source \(J_T^{\rm sol}\), \(\Lambda_{M,T}\), \(I_{M,J}\), factorisation uniforme et borne portant uniquement sur les données. \\
\(\mathfrak G_{\rm gau}\) &
\(\operatorname{Cell}_{\rm TPK},Q_{\rm inc},B,W_c\), quatre directions, six cellules, huit faces, espérances et terminal commun. &
Interface de jauge physique, semi-groupe requis, secteur physique, Yang--Mills et saut de masse. \\
\(\mathfrak G_{\rm coh}\) &
Ports, \(\kappa,\delta,C\), construction bar/coend, retenue, survie et cône d'holonomie. &
\(R_{\rm Hdg}\), \(\mathrm{ch}_{\rm loc}\), classe algébrique et expression de Hodge. \\
\(\mathfrak G_{\rm det}\) &
\(\Psi,K,z_\pm,h_*\), AW, droite déterminant, homologie, Smith et Bockstein, conservés sans exiger l'acyclicité en entrée. &
Kato--FERS, PT, comparateur local--\newline global, Tamagawa et expression BSD. \\
\bottomrule
\end{tabularx}
\end{center}

Les objets de la troisième colonne sont propres à la cible : ils ne doivent pas être réintroduits rétrospectivement comme entrées de PDF1. PDF2 les construit sur l'objet structurel complet de la deuxième colonne et conserve ses treize coordonnées, son terminal et sa généalogie.

\endgroup

\section{Typage unitaire de l'objet consommé}

Le bloc figé précédent ne s'ajoute pas à un constructeur graphique à cinq
composantes. Il le remplace. Pour chaque profondeur \(k\), soit
\(\mathcal E_k\) l'état APP--TRIT--TPK enrichi de toutes ses émissions
survivantes, et soit \(\mathcal Z_k\) le codomaine de l'unique action
corrélative. La flèche directrice est
\begin{equation}\label{eq:pdf2-constructor-unico-operaciones}
 \mathcal O_k:\mathcal E_k\longrightarrow\mathcal Z_k,
 \qquad
 \widehat x_k\longmapsto
 \mathcal O\bigl(\{\mathfrak e_\alpha:
 \alpha\in\operatorname{Ch}_k(\widehat x_k)\}\bigr).
\end{equation}
Sa valeur conserve simultanément les coordonnées
\(\mathrm{sp},\mathrm{sol},\mathrm{gau},\mathrm{coh},\mathrm{det}\)
du même état et du même ledger. Ce n'est pas un tuple de facteurs
indépendants ni le graphe tautologique de cinq applications préalablement séparées.

Les troncatures de l'état et de l'action s'écrivent
\[
 u^{\mathfrak e}_{\ell,k}:\mathcal E_\ell\to\mathcal E_k,
 \qquad
 u^{\mathcal O}_{\ell,k}:\mathcal Z_\ell\to\mathcal Z_k,
 \qquad \ell\ge k,
\]
et satisfont une seule naturalité,
\begin{equation}\label{eq:pdf2-naturalidad-accion-unica}
 u^{\mathcal O}_{\ell,k}\circ\mathcal O_\ell
 =\mathcal O_k\circ u^{\mathfrak e}_{\ell,k}.
\end{equation}
L'action passe donc à la limite sans sélectionner une histoire :
\[
 \mathcal O_\infty:
 \mathcal C_{\rm cont}^{\rm disc}:=\varprojlim_k\mathcal E_k
 \longrightarrow
 \mathcal Z_\infty:=\varprojlim_k\mathcal Z_k.
\]
Les notations utilisées dans les cinq chapitres sont des vues intrinsèques de cette
unique valeur :
\begin{equation}\label{eq:pdf2-operacion-infinita-tipificada}
 \mathfrak G_T^{\rm int}
 :=\operatorname{View}_T\!\left(
 \mathcal O_\infty(\mathcal C_{\rm cont}^{\rm disc})\right),
 \qquad
 T\in\{\mathrm{sp},\mathrm{sol},\mathrm{gau},\mathrm{coh},\mathrm{det}\}.
\end{equation}
Ici, \(\operatorname{View}_T\) expose uniquement les coordonnées requises par une
application ; elle ne crée pas de branche, de facteur, de sous-domaine ni de sélecteur. La
structure de départ de chaque résolution est toujours la valeur complète de
\(\mathcal O_\infty\), avec les quatre autres caractéristiques présentes.

\begin{falsifier}[Ordre causal et non-disgrégation]
Le raccordement échoue si l'on utilise \(\mathbb R\) ou le continu comme entrée de
\(\mathcal O_k\), si l'on remplace \(\mathcal O_k\) par cinq constructeurs
indépendants, si l'on sélectionne une seule histoire terminale ou si une vue
\(\operatorname{View}_T\) est confondue avec l'objet complet. Dans n'importe lequel
de ces cas, la conclusion du chapitre d'application cesse de se transférer.
\end{falsifier}
